
\section{Literature Review}
\label{sec:literature-review}

\subsection{Positional Bias in Long-Context Language Models}

The ability to accept a long input sequence does not necessarily imply that a
language model can use information equally well at every position. Liu et al.
(2024)~\cite{liu2024lost} investigated this distinction using multi-document
question answering and key--value retrieval. By moving relevant evidence to
different positions while retaining the underlying task, they observed that
many models performed better when the evidence appeared near the beginning or
end of the context and worse when it appeared in the middle. This pattern,
commonly called the \emph{lost-in-the-middle} effect, establishes information
position as an experimental variable rather than an incidental property of
prompt formatting.

Subsequent work has investigated ways to mitigate positional bias. Hsieh et al.
(2024)~\cite{hsieh2024found} proposed calibrating positional attention bias to
improve long-context utilization. Zhang et al. (2024)~\cite{zhang2024found}
introduced multi-scale positional encoding (Ms-PoE), which modifies positional
encoding at inference time to improve access to information in middle
positions. These studies suggest that positional performance differences can be
influenced by model mechanisms, but improvements under particular benchmarks do
not establish that the effect has disappeared across models, tasks, and context
lengths. The present project therefore treats positional bias as a measurable,
model-dependent phenomenon rather than assuming a universal U-shaped curve.

\subsection{Context Length and Effective Utilization}

A model's advertised context window specifies how much input it can accept, not
the amount of information it can reliably use. Bai et al.
(2024)~\cite{bai2024longbench} introduced LongBench, a benchmark evaluating
long-context understanding across multiple tasks, including question answering,
summarization, and code completion. Its results illustrate why performance
should be assessed on tasks requiring actual use of long inputs, rather than
inferred from supported token counts alone.

Hsieh et al. (2024)~\cite{hsieh2024ruler} introduced RULER, extending simple
needle-in-a-haystack evaluation with configurable sequence lengths and tasks
involving multiple needles, tracing, and aggregation. Their results showed that
strong performance on simple retrieval tests did not consistently carry over to
more demanding tasks or longer contexts.

Together, LongBench and RULER motivate separating nominal context capacity from
effective utilization. They also reveal a methodological challenge: when input
length increases, the amount of irrelevant material and the complexity of the
retrieval task may increase at the same time. To isolate the interaction
between length and position, the proposed study will vary total context length
and relevant-evidence position explicitly, while keeping the target answer and
task format fixed wherever feasible.

\subsection{Information Interference and Distractors}

Long contexts may contain not only irrelevant text but also competing
information that resembles the evidence needed to answer a question. This
distinction matters because a model can succeed at locating an exact string yet
struggle when it must distinguish semantically related candidates.

Modarressi et al. (2025)~\cite{modarressi2025nolima} introduced NoLiMa,
addressing a limitation of conventional needle-in-a-haystack evaluations:
lexical overlap between a question and its target passage can make retrieval
unusually easy. By reducing direct lexical matches, NoLiMa requires models to
identify less explicit associations and reports substantial performance
degradation for many models as context length grows.

However, NoLiMa does not, by itself, isolate the causal effect of
\emph{semantically similar distractors}. Removing lexical overlap between a
question and the answer passage is different from systematically increasing the
similarity of incorrect passages to the correct one. Likewise, the
multi-document question-answering experiments of Liu et al.
(2024)~\cite{liu2024lost} include non-answer passages but do not fully resolve
how distractor similarity interacts with position and length.

These distinctions motivate the third research question. A controlled
comparison should hold the correct evidence constant and vary the number and
semantic similarity of distractors, examining whether the impact differs when
relevant information appears near the beginning, middle, or end of the context.

\subsection{Synthesis and Research Gap}

The reviewed literature establishes three complementary observations.

First, evidence position can influence retrieval and question-answering
accuracy~\cite{liu2024lost}. Second, longer accepted inputs do not guarantee
consistently effective information use, particularly as task demands
increase~\cite{bai2024longbench,hsieh2024ruler}. Third, evaluation outcomes
depend on how easily relevant evidence can be distinguished from surrounding
text~\cite{modarressi2025nolima}. Methods designed to reduce positional
bias~\cite{hsieh2024found,zhang2024found} further indicate that these effects
depend on the model and its inference mechanisms.

Although existing studies investigate important aspects of long-context
utilization, the reviewed evidence does not fully characterize how evidence
position, context length, and distractor semantic similarity jointly influence
performance across contemporary language models.

The broader objective of this research is to investigate task quality ($Q$) as
a function of evidence position ($p$), total context length ($L$), and
distractor semantic similarity ($d$):

\begin{equation}
    Q = f(p,L,d).
    \label{eq:context-utilization}
\end{equation}

The research adopts a phased experimental approach. The primary investigation
focuses on the relationship between evidence position and context length, while
maintaining a controlled baseline distractor condition.

This initial phase examines whether positional effects persist in modern
language models and how their magnitude changes as context length increases.

A subsequent phase may extend the investigation to distractor semantic
similarity, depending on available time and computational resources. This
extension would examine whether semantically similar distractors alter the
relationship between evidence position, context length, and task performance.

Throughout the experiments, model parameters will remain fixed, and results
will be reported separately by model and experimental condition.
